\documentclass[11pt, a4paper]{article}
% ── Shared preamble for all RTOS lab documents ────────────────────────────────
% Usage: % ── Shared preamble for all RTOS lab documents ────────────────────────────────
% Usage: % ── Shared preamble for all RTOS lab documents ────────────────────────────────
% Usage: \input{../preamble}  (from each labNN/ directory)
% ──────────────────────────────────────────────────────────────────────────────

% ── Packages ──────────────────────────────────────────────────────────────────
\usepackage[a4paper, top=2.5cm, bottom=2.5cm, left=2.5cm, right=2.5cm, headheight=28pt]{geometry}
\usepackage[utf8]{inputenc}
\usepackage[T1]{fontenc}
\usepackage{lmodern}
\usepackage{booktabs}
\usepackage{array}
\usepackage{tabularx}
\usepackage{longtable}
\usepackage{multirow}
\usepackage{colortbl}
\usepackage{xcolor}
\usepackage{titlesec}
\usepackage{enumitem}
\usepackage{hyperref}
\usepackage{fancyhdr}
\usepackage{parskip}
\usepackage{listings}
\usepackage[most]{tcolorbox}
\usepackage{fontawesome5}
\usepackage{microtype}
\usepackage{graphicx}
\usepackage{amsmath}

% ── Colors (KMUTNB brand) ──────────────────────────────────────────────────────
\definecolor{kmutnbBlue}{RGB}{0,56,116}
\definecolor{kmutnbGold}{RGB}{186,146,36}
\definecolor{lightGray}{RGB}{245,245,245}
\definecolor{midGray}{RGB}{200,200,200}
\definecolor{codeGray}{RGB}{248,248,248}
\definecolor{codeFrame}{RGB}{220,220,220}
\definecolor{tipteal}{RGB}{0,105,92}
\definecolor{warnAmber}{RGB}{121,85,0}
\definecolor{checkGreen}{RGB}{27,94,32}

% ── Section formatting ─────────────────────────────────────────────────────────
\titleformat{\section}
  {\large\bfseries\color{kmutnbBlue}}
  {\thesection.}{0.5em}{}
  [\color{kmutnbGold}\titlerule]

\titleformat{\subsection}
  {\normalsize\bfseries\color{kmutnbBlue}}
  {\thesubsection.}{0.5em}{}

\titleformat{\subsubsection}
  {\small\bfseries\color{kmutnbBlue}}
  {\thesubsubsection.}{0.5em}{}

\titlespacing*{\section}{0pt}{1.4ex plus .2ex}{0.6ex plus .1ex}
\titlespacing*{\subsection}{0pt}{1.0ex plus .2ex}{0.4ex plus .1ex}

% ── Listings (C and bash) ──────────────────────────────────────────────────────
\lstdefinestyle{cstyle}{
  language        = C,
  basicstyle      = \ttfamily\small,
  keywordstyle    = \bfseries\color{kmutnbBlue},
  commentstyle    = \itshape\color{gray},
  stringstyle     = \color{tipteal},
  numberstyle     = \tiny\color{midGray},
  numbers         = left,
  numbersep       = 6pt,
  stepnumber      = 1,
  backgroundcolor = \color{codeGray},
  frame           = single,
  framerule       = 0.4pt,
  rulecolor       = \color{codeFrame},
  breaklines      = true,
  showstringspaces= false,
  tabsize         = 4,
  xleftmargin     = 1.5em,
  xrightmargin    = 0.5em,
  aboveskip       = 0.8em,
  belowskip       = 0.8em,
}

\lstdefinestyle{bashstyle}{
  language        = bash,
  basicstyle      = \ttfamily\small,
  keywordstyle    = \color{kmutnbBlue},
  commentstyle    = \itshape\color{gray},
  backgroundcolor = \color{black!5},
  frame           = single,
  framerule       = 0.4pt,
  rulecolor       = \color{codeFrame},
  breaklines      = true,
  showstringspaces= false,
  xleftmargin     = 1.5em,
  xrightmargin    = 0.5em,
  aboveskip       = 0.6em,
  belowskip       = 0.6em,
}

\lstdefinestyle{textstyle}{
  basicstyle      = \ttfamily\footnotesize,
  backgroundcolor = \color{black!4},
  frame           = single,
  framerule       = 0.4pt,
  rulecolor       = \color{codeFrame},
  breaklines      = true,
  xleftmargin     = 1.5em,
  xrightmargin    = 0.5em,
  aboveskip       = 0.6em,
  belowskip       = 0.6em,
}

% ── Callout boxes ──────────────────────────────────────────────────────────────
\tcbuselibrary{skins, breakable}

\newtcolorbox{tipbox}[1][]{
  enhanced, breakable,
  colback=tipteal!8, colframe=tipteal, coltitle=white,
  fonttitle=\bfseries\small, title={\faLightbulb\quad Tip#1},
  left=6pt, right=6pt, top=4pt, bottom=4pt,
  boxrule=0.8pt, arc=3pt,
}

\newtcolorbox{warnbox}[1][]{
  enhanced, breakable,
  colback=warnAmber!8, colframe=kmutnbGold, coltitle=white,
  fonttitle=\bfseries\small, title={\faExclamationTriangle\quad Note#1},
  left=6pt, right=6pt, top=4pt, bottom=4pt,
  boxrule=0.8pt, arc=3pt,
}

\newtcolorbox{checkpointbox}[1][]{
  enhanced,
  colback=kmutnbBlue!6, colframe=kmutnbBlue, coltitle=white,
  fonttitle=\bfseries\small, title={\faCheckCircle\quad Checkpoint#1},
  left=6pt, right=6pt, top=4pt, bottom=4pt,
  boxrule=0.8pt, arc=3pt,
}

\newtcolorbox{questionbox}[1][]{
  enhanced, breakable,
  colback=kmutnbGold!10, colframe=kmutnbGold, coltitle=white,
  fonttitle=\bfseries\small, title={\faQuestion\quad Question#1},
  left=6pt, right=6pt, top=4pt, bottom=4pt,
  boxrule=0.8pt, arc=3pt,
}

% ── Checkbox list ──────────────────────────────────────────────────────────────
\newlist{checklist}{itemize}{1}
\setlist[checklist]{label=\faSquare[regular], leftmargin=1.5em, noitemsep}

% ── Misc helpers ───────────────────────────────────────────────────────────────
\newcommand{\file}[1]{\texttt{\small #1}}
\newcommand{\api}[1]{\texttt{\small #1}}

% ── Title block macro ─────────────────────────────────────────────────────────
% Usage: \labtitle{03}{Aperiodic Servers \& Producer--Consumer}{2}{CLO 1 \& CLO 2}
\newcommand{\labtitle}[4]{%
  \begin{center}
    {\color{kmutnbBlue}\rule{\linewidth}{2pt}}\\[6pt]
    {\LARGE\bfseries\color{kmutnbBlue} Laboratory #1}\\[4pt]
    {\large\bfseries #2}\\[2pt]
    {\normalsize Real-Time Operating Systems \enspace\textbullet\enspace
     M.Eng.\ ECE \enspace\textbullet\enspace KMUTNB}\\[8pt]
    {\color{kmutnbGold}\rule{\linewidth}{1pt}}
  \end{center}
  \vspace{0.3cm}
  \begin{tabular}{@{}p{3.8cm} p{10cm}@{}}
    \textbf{Platform}       & NXP FRDM-MCXN236 (ARM Cortex-M33 @ 150\,MHz) \\[2pt]
    \textbf{RTOS}           & FreeRTOS v10.6 \\[2pt]
    \textbf{Toolchain}      & ARM GNU Toolchain (\texttt{arm-none-eabi-gcc}) + Make \\[2pt]
    \textbf{Estimated time} & #3 hours \\[2pt]
    \textbf{CLO mapping}    & #4 \\
  \end{tabular}
  \vspace{0.4cm}
}
  (from each labNN/ directory)
% ──────────────────────────────────────────────────────────────────────────────

% ── Packages ──────────────────────────────────────────────────────────────────
\usepackage[a4paper, top=2.5cm, bottom=2.5cm, left=2.5cm, right=2.5cm, headheight=28pt]{geometry}
\usepackage[utf8]{inputenc}
\usepackage[T1]{fontenc}
\usepackage{lmodern}
\usepackage{booktabs}
\usepackage{array}
\usepackage{tabularx}
\usepackage{longtable}
\usepackage{multirow}
\usepackage{colortbl}
\usepackage{xcolor}
\usepackage{titlesec}
\usepackage{enumitem}
\usepackage{hyperref}
\usepackage{fancyhdr}
\usepackage{parskip}
\usepackage{listings}
\usepackage[most]{tcolorbox}
\usepackage{fontawesome5}
\usepackage{microtype}
\usepackage{graphicx}
\usepackage{amsmath}

% ── Colors (KMUTNB brand) ──────────────────────────────────────────────────────
\definecolor{kmutnbBlue}{RGB}{0,56,116}
\definecolor{kmutnbGold}{RGB}{186,146,36}
\definecolor{lightGray}{RGB}{245,245,245}
\definecolor{midGray}{RGB}{200,200,200}
\definecolor{codeGray}{RGB}{248,248,248}
\definecolor{codeFrame}{RGB}{220,220,220}
\definecolor{tipteal}{RGB}{0,105,92}
\definecolor{warnAmber}{RGB}{121,85,0}
\definecolor{checkGreen}{RGB}{27,94,32}

% ── Section formatting ─────────────────────────────────────────────────────────
\titleformat{\section}
  {\large\bfseries\color{kmutnbBlue}}
  {\thesection.}{0.5em}{}
  [\color{kmutnbGold}\titlerule]

\titleformat{\subsection}
  {\normalsize\bfseries\color{kmutnbBlue}}
  {\thesubsection.}{0.5em}{}

\titleformat{\subsubsection}
  {\small\bfseries\color{kmutnbBlue}}
  {\thesubsubsection.}{0.5em}{}

\titlespacing*{\section}{0pt}{1.4ex plus .2ex}{0.6ex plus .1ex}
\titlespacing*{\subsection}{0pt}{1.0ex plus .2ex}{0.4ex plus .1ex}

% ── Listings (C and bash) ──────────────────────────────────────────────────────
\lstdefinestyle{cstyle}{
  language        = C,
  basicstyle      = \ttfamily\small,
  keywordstyle    = \bfseries\color{kmutnbBlue},
  commentstyle    = \itshape\color{gray},
  stringstyle     = \color{tipteal},
  numberstyle     = \tiny\color{midGray},
  numbers         = left,
  numbersep       = 6pt,
  stepnumber      = 1,
  backgroundcolor = \color{codeGray},
  frame           = single,
  framerule       = 0.4pt,
  rulecolor       = \color{codeFrame},
  breaklines      = true,
  showstringspaces= false,
  tabsize         = 4,
  xleftmargin     = 1.5em,
  xrightmargin    = 0.5em,
  aboveskip       = 0.8em,
  belowskip       = 0.8em,
}

\lstdefinestyle{bashstyle}{
  language        = bash,
  basicstyle      = \ttfamily\small,
  keywordstyle    = \color{kmutnbBlue},
  commentstyle    = \itshape\color{gray},
  backgroundcolor = \color{black!5},
  frame           = single,
  framerule       = 0.4pt,
  rulecolor       = \color{codeFrame},
  breaklines      = true,
  showstringspaces= false,
  xleftmargin     = 1.5em,
  xrightmargin    = 0.5em,
  aboveskip       = 0.6em,
  belowskip       = 0.6em,
}

\lstdefinestyle{textstyle}{
  basicstyle      = \ttfamily\footnotesize,
  backgroundcolor = \color{black!4},
  frame           = single,
  framerule       = 0.4pt,
  rulecolor       = \color{codeFrame},
  breaklines      = true,
  xleftmargin     = 1.5em,
  xrightmargin    = 0.5em,
  aboveskip       = 0.6em,
  belowskip       = 0.6em,
}

% ── Callout boxes ──────────────────────────────────────────────────────────────
\tcbuselibrary{skins, breakable}

\newtcolorbox{tipbox}[1][]{
  enhanced, breakable,
  colback=tipteal!8, colframe=tipteal, coltitle=white,
  fonttitle=\bfseries\small, title={\faLightbulb\quad Tip#1},
  left=6pt, right=6pt, top=4pt, bottom=4pt,
  boxrule=0.8pt, arc=3pt,
}

\newtcolorbox{warnbox}[1][]{
  enhanced, breakable,
  colback=warnAmber!8, colframe=kmutnbGold, coltitle=white,
  fonttitle=\bfseries\small, title={\faExclamationTriangle\quad Note#1},
  left=6pt, right=6pt, top=4pt, bottom=4pt,
  boxrule=0.8pt, arc=3pt,
}

\newtcolorbox{checkpointbox}[1][]{
  enhanced,
  colback=kmutnbBlue!6, colframe=kmutnbBlue, coltitle=white,
  fonttitle=\bfseries\small, title={\faCheckCircle\quad Checkpoint#1},
  left=6pt, right=6pt, top=4pt, bottom=4pt,
  boxrule=0.8pt, arc=3pt,
}

\newtcolorbox{questionbox}[1][]{
  enhanced, breakable,
  colback=kmutnbGold!10, colframe=kmutnbGold, coltitle=white,
  fonttitle=\bfseries\small, title={\faQuestion\quad Question#1},
  left=6pt, right=6pt, top=4pt, bottom=4pt,
  boxrule=0.8pt, arc=3pt,
}

% ── Checkbox list ──────────────────────────────────────────────────────────────
\newlist{checklist}{itemize}{1}
\setlist[checklist]{label=\faSquare[regular], leftmargin=1.5em, noitemsep}

% ── Misc helpers ───────────────────────────────────────────────────────────────
\newcommand{\file}[1]{\texttt{\small #1}}
\newcommand{\api}[1]{\texttt{\small #1}}

% ── Title block macro ─────────────────────────────────────────────────────────
% Usage: \labtitle{03}{Aperiodic Servers \& Producer--Consumer}{2}{CLO 1 \& CLO 2}
\newcommand{\labtitle}[4]{%
  \begin{center}
    {\color{kmutnbBlue}\rule{\linewidth}{2pt}}\\[6pt]
    {\LARGE\bfseries\color{kmutnbBlue} Laboratory #1}\\[4pt]
    {\large\bfseries #2}\\[2pt]
    {\normalsize Real-Time Operating Systems \enspace\textbullet\enspace
     M.Eng.\ ECE \enspace\textbullet\enspace KMUTNB}\\[8pt]
    {\color{kmutnbGold}\rule{\linewidth}{1pt}}
  \end{center}
  \vspace{0.3cm}
  \begin{tabular}{@{}p{3.8cm} p{10cm}@{}}
    \textbf{Platform}       & NXP FRDM-MCXN236 (ARM Cortex-M33 @ 150\,MHz) \\[2pt]
    \textbf{RTOS}           & FreeRTOS v10.6 \\[2pt]
    \textbf{Toolchain}      & ARM GNU Toolchain (\texttt{arm-none-eabi-gcc}) + Make \\[2pt]
    \textbf{Estimated time} & #3 hours \\[2pt]
    \textbf{CLO mapping}    & #4 \\
  \end{tabular}
  \vspace{0.4cm}
}
  (from each labNN/ directory)
% ──────────────────────────────────────────────────────────────────────────────

% ── Packages ──────────────────────────────────────────────────────────────────
\usepackage[a4paper, top=2.5cm, bottom=2.5cm, left=2.5cm, right=2.5cm, headheight=28pt]{geometry}
\usepackage[utf8]{inputenc}
\usepackage[T1]{fontenc}
\usepackage{lmodern}
\usepackage{booktabs}
\usepackage{array}
\usepackage{tabularx}
\usepackage{longtable}
\usepackage{multirow}
\usepackage{colortbl}
\usepackage{xcolor}
\usepackage{titlesec}
\usepackage{enumitem}
\usepackage{hyperref}
\usepackage{fancyhdr}
\usepackage{parskip}
\usepackage{listings}
\usepackage[most]{tcolorbox}
\usepackage{fontawesome5}
\usepackage{microtype}
\usepackage{graphicx}
\usepackage{amsmath}

% ── Colors (KMUTNB brand) ──────────────────────────────────────────────────────
\definecolor{kmutnbBlue}{RGB}{0,56,116}
\definecolor{kmutnbGold}{RGB}{186,146,36}
\definecolor{lightGray}{RGB}{245,245,245}
\definecolor{midGray}{RGB}{200,200,200}
\definecolor{codeGray}{RGB}{248,248,248}
\definecolor{codeFrame}{RGB}{220,220,220}
\definecolor{tipteal}{RGB}{0,105,92}
\definecolor{warnAmber}{RGB}{121,85,0}
\definecolor{checkGreen}{RGB}{27,94,32}

% ── Section formatting ─────────────────────────────────────────────────────────
\titleformat{\section}
  {\large\bfseries\color{kmutnbBlue}}
  {\thesection.}{0.5em}{}
  [\color{kmutnbGold}\titlerule]

\titleformat{\subsection}
  {\normalsize\bfseries\color{kmutnbBlue}}
  {\thesubsection.}{0.5em}{}

\titleformat{\subsubsection}
  {\small\bfseries\color{kmutnbBlue}}
  {\thesubsubsection.}{0.5em}{}

\titlespacing*{\section}{0pt}{1.4ex plus .2ex}{0.6ex plus .1ex}
\titlespacing*{\subsection}{0pt}{1.0ex plus .2ex}{0.4ex plus .1ex}

% ── Listings (C and bash) ──────────────────────────────────────────────────────
\lstdefinestyle{cstyle}{
  language        = C,
  basicstyle      = \ttfamily\small,
  keywordstyle    = \bfseries\color{kmutnbBlue},
  commentstyle    = \itshape\color{gray},
  stringstyle     = \color{tipteal},
  numberstyle     = \tiny\color{midGray},
  numbers         = left,
  numbersep       = 6pt,
  stepnumber      = 1,
  backgroundcolor = \color{codeGray},
  frame           = single,
  framerule       = 0.4pt,
  rulecolor       = \color{codeFrame},
  breaklines      = true,
  showstringspaces= false,
  tabsize         = 4,
  xleftmargin     = 1.5em,
  xrightmargin    = 0.5em,
  aboveskip       = 0.8em,
  belowskip       = 0.8em,
}

\lstdefinestyle{bashstyle}{
  language        = bash,
  basicstyle      = \ttfamily\small,
  keywordstyle    = \color{kmutnbBlue},
  commentstyle    = \itshape\color{gray},
  backgroundcolor = \color{black!5},
  frame           = single,
  framerule       = 0.4pt,
  rulecolor       = \color{codeFrame},
  breaklines      = true,
  showstringspaces= false,
  xleftmargin     = 1.5em,
  xrightmargin    = 0.5em,
  aboveskip       = 0.6em,
  belowskip       = 0.6em,
}

\lstdefinestyle{textstyle}{
  basicstyle      = \ttfamily\footnotesize,
  backgroundcolor = \color{black!4},
  frame           = single,
  framerule       = 0.4pt,
  rulecolor       = \color{codeFrame},
  breaklines      = true,
  xleftmargin     = 1.5em,
  xrightmargin    = 0.5em,
  aboveskip       = 0.6em,
  belowskip       = 0.6em,
}

% ── Callout boxes ──────────────────────────────────────────────────────────────
\tcbuselibrary{skins, breakable}

\newtcolorbox{tipbox}[1][]{
  enhanced, breakable,
  colback=tipteal!8, colframe=tipteal, coltitle=white,
  fonttitle=\bfseries\small, title={\faLightbulb\quad Tip#1},
  left=6pt, right=6pt, top=4pt, bottom=4pt,
  boxrule=0.8pt, arc=3pt,
}

\newtcolorbox{warnbox}[1][]{
  enhanced, breakable,
  colback=warnAmber!8, colframe=kmutnbGold, coltitle=white,
  fonttitle=\bfseries\small, title={\faExclamationTriangle\quad Note#1},
  left=6pt, right=6pt, top=4pt, bottom=4pt,
  boxrule=0.8pt, arc=3pt,
}

\newtcolorbox{checkpointbox}[1][]{
  enhanced,
  colback=kmutnbBlue!6, colframe=kmutnbBlue, coltitle=white,
  fonttitle=\bfseries\small, title={\faCheckCircle\quad Checkpoint#1},
  left=6pt, right=6pt, top=4pt, bottom=4pt,
  boxrule=0.8pt, arc=3pt,
}

\newtcolorbox{questionbox}[1][]{
  enhanced, breakable,
  colback=kmutnbGold!10, colframe=kmutnbGold, coltitle=white,
  fonttitle=\bfseries\small, title={\faQuestion\quad Question#1},
  left=6pt, right=6pt, top=4pt, bottom=4pt,
  boxrule=0.8pt, arc=3pt,
}

% ── Checkbox list ──────────────────────────────────────────────────────────────
\newlist{checklist}{itemize}{1}
\setlist[checklist]{label=\faSquare[regular], leftmargin=1.5em, noitemsep}

% ── Misc helpers ───────────────────────────────────────────────────────────────
\newcommand{\file}[1]{\texttt{\small #1}}
\newcommand{\api}[1]{\texttt{\small #1}}

% ── Title block macro ─────────────────────────────────────────────────────────
% Usage: \labtitle{03}{Aperiodic Servers \& Producer--Consumer}{2}{CLO 1 \& CLO 2}
\newcommand{\labtitle}[4]{%
  \begin{center}
    {\color{kmutnbBlue}\rule{\linewidth}{2pt}}\\[6pt]
    {\LARGE\bfseries\color{kmutnbBlue} Laboratory #1}\\[4pt]
    {\large\bfseries #2}\\[2pt]
    {\normalsize Real-Time Operating Systems \enspace\textbullet\enspace
     M.Eng.\ ECE \enspace\textbullet\enspace KMUTNB}\\[8pt]
    {\color{kmutnbGold}\rule{\linewidth}{1pt}}
  \end{center}
  \vspace{0.3cm}
  \begin{tabular}{@{}p{3.8cm} p{10cm}@{}}
    \textbf{Platform}       & NXP FRDM-MCXN236 (ARM Cortex-M33 @ 150\,MHz) \\[2pt]
    \textbf{RTOS}           & FreeRTOS v10.6 \\[2pt]
    \textbf{Toolchain}      & ARM GNU Toolchain (\texttt{arm-none-eabi-gcc}) + Make \\[2pt]
    \textbf{Estimated time} & #3 hours \\[2pt]
    \textbf{CLO mapping}    & #4 \\
  \end{tabular}
  \vspace{0.4cm}
}


% ── Header / Footer ────────────────────────────────────────────────────────────
\pagestyle{fancy}
\fancyhf{}
\fancyhead[L]{\small\color{kmutnbBlue}\textbf{KMUTNB -- Faculty of Engineering}}
\fancyhead[R]{\small\color{kmutnbBlue}Dept.\ of Electrical and Computer Engineering}
\fancyfoot[L]{\small\color{kmutnbBlue}Real-Time Operating Systems -- M.Eng.\ ECE}
\fancyfoot[C]{\small\thepage}
\fancyfoot[R]{\small\color{kmutnbBlue}Lab 05}
\renewcommand{\headrulewidth}{0.4pt}
\renewcommand{\headrule}{\hbox to\headwidth{%
  \color{kmutnbGold}\leaders\hrule height\headrulewidth\hfill}}
\renewcommand{\footrulewidth}{0.4pt}
\renewcommand{\footrule}{\hbox to\headwidth{%
  \color{midGray}\leaders\hrule height\footrulewidth\hfill}}

\hypersetup{
  colorlinks = true,
  linkcolor  = kmutnbBlue,
  urlcolor   = kmutnbBlue,
  citecolor  = kmutnbBlue,
  pdftitle   = {Lab 05 -- Deadlock Detection \& Watchdog Timers},
  pdfauthor  = {KMUTNB RTOS Course},
}

% ══════════════════════════════════════════════════════════════════════════════
\begin{document}

\labtitle{05}{Deadlock Detection \& Watchdog Timers}{3}{CLO 4}

% ── Objectives ────────────────────────────────────────────────────────────────
\section{Objectives}

By completing this lab you will be able to:

\begin{enumerate}[noitemsep, leftmargin=*]
  \item \textbf{Induce} a deadlock between two tasks holding two mutexes in opposite
        order.
  \item \textbf{Prevent} deadlock using the resource-ordering convention.
  \item \textbf{Configure} the MCXN236 hardware watchdog (WDOG32) and observe a reset.
  \item \textbf{Implement} a software watchdog pattern using task notifications.
  \item \textbf{Isolate} a hung task --- log its identity and trigger a controlled reset.
\end{enumerate}

% ── Prerequisites ─────────────────────────────────────────────────────────────
\section{Prerequisites}

\begin{checklist}
  \item Lab 04 complete --- mutexes and SystemView working.
  \item ARM GNU Toolchain and \texttt{pyocd} working.
\end{checklist}

% ── Background ────────────────────────────────────────────────────────────────
\section{Background}

\subsection{Deadlock and the Four Coffman Conditions}

A deadlock requires all four conditions simultaneously:

\begin{enumerate}[noitemsep, leftmargin=*]
  \item \textbf{Mutual exclusion} --- resources are non-shareable.
  \item \textbf{Hold and wait} --- a task holds at least one resource while waiting for
        another.
  \item \textbf{No preemption} --- resources are released only voluntarily.
  \item \textbf{Circular wait} --- a cycle exists in the resource-allocation graph.
\end{enumerate}

FreeRTOS mutexes satisfy conditions 1--3 by design. Your task is to prevent condition~4
using \textbf{resource ordering}: assign a global lock order and always acquire resources
in ascending order.

\subsection{Hardware Watchdog}

The MCXN236 WDOG32 generates a system reset if not refreshed within a configured window.
This is the last line of defence against a hung system.

% ── Project Structure ─────────────────────────────────────────────────────────
\section{Project Structure}

\begin{lstlisting}[style=textstyle]
lab05_watchdog/
  Makefile
  linker_MCXN236.ld
  src/
    main.c               <- task creation, WDOG init
    deadlock_tasks.c/.h  <- Parts A & B
    watchdog.c/.h        <- Parts C & D hardware + software watchdog
    FreeRTOSConfig.h
    uart.h / uart.c
  FreeRTOS-Kernel/
\end{lstlisting}

% ── Part A ─────────────────────────────────────────────────────────────────────
\section{Part A --- Induce a Deadlock}

\subsection{Step 1 --- Implement the Deadlock Scenario}

Create two mutexes and two tasks that acquire them in opposite order:

\begin{lstlisting}[style=cstyle]
/* main.c */
SemaphoreHandle_t xMutexA;
SemaphoreHandle_t xMutexB;

xMutexA = xSemaphoreCreateMutex();
xMutexB = xSemaphoreCreateMutex();
\end{lstlisting}

\begin{lstlisting}[style=cstyle]
/* deadlock_tasks.c */
void vTaskAlpha(void *pv)  /* prio 2 */
{
    for (;;) {
        vTaskDelay(pdMS_TO_TICKS(10));
        uart_printf("[A] taking MutexA...\r\n");
        xSemaphoreTake(xMutexA, portMAX_DELAY);
        uart_printf("[A] took MutexA. Taking MutexB...\r\n");
        vTaskDelay(pdMS_TO_TICKS(5));
        xSemaphoreTake(xMutexB, portMAX_DELAY);
        uart_printf("[A] took both -- doing work\r\n");
        xSemaphoreGive(xMutexB);
        xSemaphoreGive(xMutexA);
    }
}

void vTaskBeta(void *pv)   /* prio 2 */
{
    for (;;) {
        vTaskDelay(pdMS_TO_TICKS(10));
        uart_printf("[B] taking MutexB...\r\n");
        xSemaphoreTake(xMutexB, portMAX_DELAY);
        uart_printf("[B] took MutexB. Taking MutexA...\r\n");
        vTaskDelay(pdMS_TO_TICKS(5));
        xSemaphoreTake(xMutexA, portMAX_DELAY);  /* <- opposite order to Task A */
        uart_printf("[B] took both -- doing work\r\n");
        xSemaphoreGive(xMutexA);
        xSemaphoreGive(xMutexB);
    }
}
\end{lstlisting}

\subsection{Step 2 --- Build, Flash, and Observe}

\begin{lstlisting}[style=bashstyle]
make setup && make && make flash
\end{lstlisting}

Watch the terminal. Within a few seconds, both \texttt{[A]} and \texttt{[B]} will stop
printing ``took both'' --- the system is deadlocked.

\begin{checkpointbox}[~A]
  UART capture showing the last lines before deadlock, with both tasks blocked.
\end{checkpointbox}

\begin{questionbox}[~A1]
  Draw the resource allocation graph at the moment of deadlock. Label each edge as an
  assignment or a request edge. Which Coffman condition does the circular edge represent?
\end{questionbox}

\begin{questionbox}[~A2]
  The deadlock is \textbf{non-deterministic} --- it does not always happen on the first
  iteration. What determines how quickly it manifests? Would the deadlock still be
  possible if both tasks ran at different priorities?
\end{questionbox}

% ── Part B ─────────────────────────────────────────────────────────────────────
\section{Part B --- Prevent with Resource Ordering}

\subsection{Step 1 --- Apply Resource Ordering}

Assign a global lock order constant and enforce it in both tasks:

\begin{lstlisting}[style=cstyle]
/* deadlock_tasks.c -- FIXED */
#define LOCK_ORDER_MUTEX_A  1
#define LOCK_ORDER_MUTEX_B  2
/* Rule: always acquire in ascending order */

/* vTaskAlpha unchanged -- already acquires A then B (correct order) */

void vTaskBeta(void *pv)   /* FIXED: acquire A before B */
{
    for (;;) {
        vTaskDelay(pdMS_TO_TICKS(10));
        uart_printf("[B] taking MutexA (order fix)...\r\n");
        xSemaphoreTake(xMutexA, portMAX_DELAY);   /* A first */
        uart_printf("[B] taking MutexB...\r\n");
        xSemaphoreTake(xMutexB, portMAX_DELAY);   /* B second */
        uart_printf("[B] took both -- doing work\r\n");
        xSemaphoreGive(xMutexB);
        xSemaphoreGive(xMutexA);
    }
}
\end{lstlisting}

\subsection{Step 2 --- Verify No Deadlock}

Run for at least 60 seconds. Both tasks should continue printing ``took both''
indefinitely.

\begin{checkpointbox}[~B]
  UART capture of 10+ seconds of normal operation with no deadlock.
\end{checkpointbox}

\begin{questionbox}[~B1]
  Resource ordering prevents deadlock by breaking which Coffman condition? Explain
  precisely why a total order on acquisition prevents a cycle in the
  resource-allocation graph.
\end{questionbox}

\begin{questionbox}[~B2]
  Your system has 6 mutexes protecting 6 different subsystems. How do you choose the
  lock order? What property must the ordering have? Is it always possible?
\end{questionbox}

% ── Part C ─────────────────────────────────────────────────────────────────────
\section{Part C --- Hardware Watchdog}

\subsection{Step 1 --- Configure WDOG32}

The MCXN236 WDOG32 resets the chip if not refreshed within a timeout window.

\begin{lstlisting}[style=cstyle]
/* watchdog.c */
#include "fsl_wdog32.h"

void WDOG_Init_1s(void)
{
    wdog32_config_t cfg;
    WDOG32_GetDefaultConfig(&cfg);
    cfg.timeoutValue = 0x100U;   /* ~1 second at LPO 128 Hz */
    cfg.enableWdog32 = true;
    cfg.clockSource  = kWDOG32_ClockSourceLpoClk;
    cfg.prescaler    = kWDOG32_ClockPrescalerDivide1;
    WDOG32_Init(WDOG0, &cfg);
}

void WDOG_Kick(void)
{
    WDOG32_Refresh(WDOG0);
}
\end{lstlisting}

\subsection{Step 2 --- Normal Operation}

Create a high-priority watchdog task that kicks every 500\,ms:

\begin{lstlisting}[style=cstyle]
void vWatchdogTask(void *pv)  /* prio 4 -- highest */
{
    WDOG_Init_1s();
    for (;;) {
        vTaskDelay(pdMS_TO_TICKS(500));
        WDOG_Kick();
        uart_printf("[WDG] kicked\r\n");
    }
}
\end{lstlisting}

\subsection{Step 3 --- Trigger the Watchdog Reset}

Add a \texttt{\#define SIMULATE\_HANG 1} build flag. When set, \texttt{vWatchdogTask}
sleeps for 2\,seconds instead of kicking:

\begin{lstlisting}[style=cstyle]
#ifdef SIMULATE_HANG
    vTaskDelay(pdMS_TO_TICKS(2000));   /* miss deadline -> watchdog fires */
#else
    WDOG_Kick();
#endif
\end{lstlisting}

Build with \texttt{make HANG=1 \&\& make HANG=1 flash}.

In \file{main.c}, before the scheduler starts, check the reset cause:

\begin{lstlisting}[style=cstyle]
if (WDOG32_GetStatusFlags(WDOG0) & kWDOG32_InterruptFlag) {
    uart_printf("[BOOT] Reset caused by watchdog!\r\n");
}
\end{lstlisting}

\begin{checkpointbox}[~C]
  UART capture showing \texttt{[BOOT] Reset caused by watchdog!} on the reboot
  following the simulated hang.
\end{checkpointbox}

\begin{questionbox}[~C1]
  The WDOG32 uses the LPO (Low-Power Oscillator) at 128\,Hz. Calculate the exact
  timeout in milliseconds for \texttt{timeoutValue = 0x100}. Why is a low-power
  oscillator used rather than the main system clock?
\end{questionbox}

\begin{questionbox}[~C2]
  A \textbf{window watchdog} adds a minimum refresh time --- you cannot kick too early
  OR too late. How would you configure a window watchdog to catch a ``too-fast'' kick?
  Sketch the \texttt{cfg.windowValue} setting.
\end{questionbox}

% ── Part D ─────────────────────────────────────────────────────────────────────
\section{Part D --- Software Watchdog with Task Monitoring}

\subsection{Step 1 --- Multi-Task Monitoring Pattern}

The hardware watchdog is kicked by a dedicated monitor task that only kicks if
\textbf{all} worker tasks have checked in during the last interval:

\begin{lstlisting}[style=cstyle]
/* watchdog.c -- software monitor */
#define NUM_WORKERS  3
static volatile uint8_t worker_alive[NUM_WORKERS];

void wdog_checkin(uint8_t worker_id)
{
    worker_alive[worker_id] = 1;
}

void vWatchdogMonitor(void *pv)  /* prio 4 */
{
    WDOG_Init_1s();
    for (;;) {
        vTaskDelay(pdMS_TO_TICKS(500));
        bool all_ok = true;
        for (int i = 0; i < NUM_WORKERS; i++) {
            if (!worker_alive[i]) {
                uart_printf("[WDG] worker %d silent -- NOT kicking!\r\n", i);
                all_ok = false;
            }
            worker_alive[i] = 0;
        }
        if (all_ok) {
            WDOG_Kick();
            uart_printf("[WDG] all workers OK -- kicked\r\n");
        }
    }
}

void vWorkerTask(void *pv)
{
    uint8_t id = (uint8_t)(uintptr_t)pv;
    for (;;) {
        vTaskDelay(pdMS_TO_TICKS(200));
        wdog_checkin(id);
        uart_printf("[W%d] working\r\n", id);
    }
}
\end{lstlisting}

\subsection{Step 2 --- Simulate a Hanging Worker}

Add a flag to stop one worker checking in after 5 seconds:

\begin{lstlisting}[style=cstyle]
if (id == 1 && xTaskGetTickCount() > pdMS_TO_TICKS(5000)) {
    for (;;) {}   /* hang -- no checkin */
}
\end{lstlisting}

After the hang, the monitor stops kicking within 500\,ms, and the watchdog fires
$\approx$1\,second later.

\begin{checkpointbox}[~D]
  UART capture showing the monitor detecting the silent worker, then the board reset
  and reboot message.
\end{checkpointbox}

\begin{questionbox}[~D1]
  In this pattern, the monitor task itself could hang. How would you guard against
  that? (Hint: think about who monitors the monitor.)
\end{questionbox}

\begin{questionbox}[~D2]
  The worker uses a \texttt{volatile uint8\_t} flag rather than a task notification.
  What are the pros and cons of each approach? Which is safer in an RTOS context and
  why?
\end{questionbox}

% ── Deliverables ───────────────────────────────────────────────────────────────
\section{Deliverables}

\renewcommand{\arraystretch}{1.3}
\begin{tabular}{@{} c p{4.5cm} p{7cm} @{}}
  \toprule
  \textbf{\#} & \textbf{Item} & \textbf{Details} \\
  \midrule
  1 & Part A UART capture  & Deadlock state visible \\
  2 & Part A RAG diagram   & Resource-allocation graph (hand-drawn or tool) \\
  3 & Part B UART capture  & 60+ seconds of normal operation \\
  4 & Part C UART capture  & Watchdog reset cause message on reboot \\
  5 & Part D UART capture  & Monitor detecting silent worker + reset \\
  6 & Written answers      & Questions A1--A2, B1--B2, C1--C2, D1--D2 \\
  \bottomrule
\end{tabular}
\renewcommand{\arraystretch}{1}

% ── Key Concepts ───────────────────────────────────────────────────────────────
\section{Key Concepts Demonstrated}

\renewcommand{\arraystretch}{1.3}
\begin{tabular}{@{} p{7cm} p{6cm} @{}}
  \toprule
  \textbf{Concept} & \textbf{Where you saw it} \\
  \midrule
  Deadlock induction --- opposite lock order    & Part A \\
  Coffman condition 4 --- circular wait         & Part A diagram \\
  Resource ordering --- deadlock prevention     & Part B \\
  WDOG32 hardware watchdog configuration        & Part C \\
  Reset cause detection at boot                 & Part C \\
  Software watchdog with per-task monitoring    & Part D \\
  \bottomrule
\end{tabular}
\renewcommand{\arraystretch}{1}

% ── Reference ─────────────────────────────────────────────────────────────────
\section{Reference}

\renewcommand{\arraystretch}{1.3}
\begin{tabular}{@{} p{5.5cm} p{7.5cm} @{}}
  \toprule
  \textbf{Resource} & \textbf{Relevant sections} \\
  \midrule
  Coffman, Elphick \& Shoshani (1971) &
    ``System deadlocks'' (\textit{ACM Computing Surveys}) \\
  Barry, \textit{Mastering the FreeRTOS Kernel} & Ch.~6 (mutexes) \\
  NXP MCXN236 Reference Manual & Chapter WDOG32 \\
  NXP MCUXpresso SDK & \texttt{fsl\_wdog32.h}, \texttt{WDOG32\_GetDefaultConfig} \\
  \bottomrule
\end{tabular}
\renewcommand{\arraystretch}{1}

\end{document}
